\section{Continuous Assessment 1 (CA-1) --- Set B}
\label{sec:ca1_set_b}

\LPUPaperHeader{Continuous Assessment (CA-1)}{B}{60 Minutes}{30 Marks}{CO1 (Introduction \& Sustainability), CO2 (Natural Resources \& Ecosystems)}

\begin{center}
\sffamily\textbf{\color{AlertRed}Note: Attempt all six questions. Each question carries 5 marks.}
\end{center}

\questionitem{1}{
    Analyze the concept of Ecological Carrying Capacity ($K$). Explain the difference between biological carrying capacity and cultural carrying capacity, and discuss what occurs when human ecological footprint exceeds biocapacity.
}{5}{CO1}{L2 Understand}

\questionitem{2}{
    Examine the United Nations Sustainable Development Goals (SDGs). Detail the specific mandates of SDG 6 (Clean Water), SDG 7 (Affordable Clean Energy), SDG 11 (Sustainable Cities), and SDG 13 (Climate Action).
}{5}{CO1}{L2 Understand}

\questionitem{3}{
    Describe the soil profile horizons ($O, A, E, B, C, R$). Explain the biogeochemical role of humus in the A-horizon and the causes of soil salinization in canal-irrigated arid agricultural tracts.
}{5}{CO2}{L3 Apply}

\questionitem{4}{
    Evaluate the National Forest Policy (1988) targets for plain and mountainous regions in India. Discuss the ecological implications of large-scale monoculture plantations versus mixed native forests.
}{5}{CO2}{L3 Apply}

\questionitem{5}{
    Explain the difference between Gross Primary Productivity (GPP), Net Primary Productivity (NPP), and Net Ecosystem Productivity (NEP). If an agro-ecosystem has $\text{GPP} = 35{,}000\text{ kcal/m}^2/\text{yr}$ and autotrophic respiration is $45\%$, calculate NPP.
}{5}{CO2}{L3 Apply}

\questionitem{6}{
    Compare primary and secondary ecological succession. Outline the six seral stages in the Clementsian succession model (Nudation to Climax) and state three thermodynamic trends observed as succession advances toward climax.
}{5}{CO2}{L4 Analyze}

\newpage
\subsection*{Solutions and Marking Rubrics --- CA-1 Set B}

\begin{solutionbox}[Solution to Question 5 --- Productivity Calculations]
\textbf{1. Definitions (2 Marks):}\\
$\text{GPP}$: Total rate of organic chemical energy fixed by photosynthesis.\\
$\text{NPP} = \text{GPP} - R_a$: Net energy accumulated as biomass after autotrophic cellular respiration.\\
$\text{NEP} = \text{NPP} - R_h$: Net ecosystem carbon storage after heterotrophic consumer respiration.\\
\textbf{2. Calculation (2 Marks):}\\
Given: $\text{GPP} = 35{,}000\text{ kcal/m}^2/\text{yr}$, $R_a = 0.45 \times \text{GPP}$.\\
$\text{NPP} = \text{GPP} \times (1 - 0.45) = 35{,}000 \times 0.55 = 19{,}250\text{ kcal/m}^2/\text{yr}$.\\
\textbf{3. Ecological Interpretation (1 Mark):}\\
$19{,}250\text{ kcal/m}^2/\text{yr}$ represents the net chemical energy available to heterotrophic herbivores and harvestable agricultural yield.
\end{solutionbox}

\begin{solutionbox}[Solution to Question 6 --- Ecological Succession Dynamics]
\textbf{1. Primary vs Secondary (1.5 Marks):}\\
Primary initiates on sterile substrates without prior soil (bare rock, lava); secondary initiates on disturbed sites with intact soil and seed banks (wildfire, abandoned fields).\\
\textbf{2. Seral Stages (2 Marks):}\\
Nudation $\to$ Invasion (Migration) $\to$ Ecesis (Establishment) $\to$ Aggregation $\to$ Competition $\to$ Reaction $\to$ Climax Stabilization.\\
\textbf{3. Thermodynamic Trends (1.5 Marks):}\\
$P/R$ ratio approaches $1.0$; species diversity ($H'$) increases; food chains evolve into interconnected food webs.
\end{solutionbox}
\newpage
